\uuid{EABF}
\titre{Domaine de définition}
\auteur{bodin ; hanani ; mzari}
\organisation{exo7}
\licence{CC BY-NC-SA 4.0}
\datecreate{2020-01-20}
\niveau{L1}
\module{Analyse}
\chapitre{Fonctions circulaires et hyperboliques inverses}
\sousChapitre{Autre}
\difficulte{1}
\theme{Fonctions usuelles, domaine de définition}
\exo7qcmid{lille-1.fonctions-usuelles.003}

\contenu{
\begin{question}
\qclassification{126.00}
\qtags{motcle=domaine de définition}
Soit $ f(x)= \ln(\frac{2+x}{2-x}) $ et $g(x)=x^x$. On notera $D_f$ et $D_g$ le domaine de définition des fonctions $f$ et $g$ respectivement. Quelles sont les assertions vraies ?
\begin{answers}
    \bad{$D_f=\Rr\backslash\{2\}$}

    \good{$D_f=]-2,2[$}

    \bad{$D_g=\Rr$}

    \good{$D_g=]0,+\infty[$}
\end{answers}
\begin{explanations}
$f$ est définie  si $x\neq 2$ et $\frac{2+x}{2-x}\ge 0$. On déduit que
 $D_f=]-2,2[$. Par définition, $g(x)=e^{x\ln x}$. Donc $g$ est définie si $x>0$.
\end{explanations}
\end{question}
}
